\begin{ouverture}[son]{14}{Émettre et transmettre un son}{Ondes et signaux}{Observer · Modéliser · Expérimenter · Expliquer}
\objectif{Expliquer émission, propagation et réception d’un son.}
\objectif{Relier distance, durée et vitesse de propagation.}
\objectif{Concevoir une mesure avec deux microphones.}
\prerequis{Proportionnalité, conversions et lecture de graphiques.}
\end{ouverture}
\essentiel
\essSection{1}{Une vibration et un milieu}
Une corde, une membrane ou une lame qui vibre peut mettre le milieu environnant en vibration. Une caisse de résonance facilite le couplage avec l'air et peut renforcer le son émis. Le signal sonore se propage de proche en proche dans un milieu matériel ; il ne se propage pas dans le vide.
\schema[corrige]{propagation-son}
Les particules du milieu oscillent autour de leur position ; elles ne voyagent pas toutes de la source jusqu'à l'auditeur. Une onde transporte une perturbation et de l'énergie, sans transport global de matière sur cette distance.
\essSection{2}{Une vitesse de propagation}
\begin{loi}[Propagation dans un milieu uniforme]Pour une distance parcourue $d$ et une durée de propagation $\Delta t$, $v=d/\Delta t$. Dans l'air proche de la température ambiante, on utilise souvent $v\simeq\qty{340}{m.s^{-1}}$. La valeur dépend du milieu et de ses conditions.\end{loi}
Le son se propage généralement plus vite dans l'eau que dans l'air. Les valeurs nécessaires sont données dans les exercices. La vitesse de propagation du signal n'est pas la vitesse de vibration d'une particule du milieu.
\begin{exemple}[Un écho]Un obstacle renvoie un signal en 0,200 s. Le trajet est un aller-retour : $2d=v\Delta t$, donc $d=340\times0,200/2=\qty{34,0}{m}$.\end{exemple}
\essSection{3}{Mesurer un retard}
Deux microphones séparés d'une distance connue reçoivent un même signal à des dates différentes. Si la source et les microphones sont alignés dans le bon ordre, la différence de trajets vaut la distance entre les microphones.
\schema[corrige]{retard-son}
\begin{savoirfaire}[Une mesure exploitable]Repérer le même détail du signal sur les deux enregistrements ; lire le retard et non la durée totale de l'acquisition ; vérifier les unités de temps. Augmenter raisonnablement la distance peut rendre le retard plus facile à mesurer.\end{savoirfaire}
\essSection{4}{Expérimenter et modéliser}
\begin{experience}[Mesurer la vitesse du son]Aligner une source impulsionnelle et deux microphones. Mesurer leur séparation, enregistrer simultanément, lire le retard et calculer $v$. Refaire à plusieurs distances, discuter les échos de la salle et la lecture temporelle. Employer un signal bref de niveau raisonnable ; ne pas produire de choc sonore près des oreilles.\end{experience}
\exercices
\rubrique{Prendre confiance}
\begin{blocLivre}{ex-14-1}\begin{exercice}[1]{Une source sonore}
Une corde de guitare est pincée.
\begin{questions}
\item Quel objet vibre à l’origine du son ?
\item Quel intérêt présente la caisse de la guitare ?
\end{questions}
\end{exercice}\end{blocLivre}
\begin{blocLivre}{ex-14-2}\begin{exercice}[1]{Le rôle du milieu}
Une sonnerie est placée dans une enceinte dont on retire progressivement l’air.
\schema{propagation-son}
\begin{questions}
\item Pourquoi le son transmis à travers le gaz s’affaiblit-il ?
\item Un son pourrait-il traverser un vide parfait ?
\end{questions}
\end{exercice}\end{blocLivre}
\begin{blocLivre}{ex-14-3}\begin{exercice}[1]{Une durée de propagation}
Un signal parcourt 68,0 m dans l’air à 340 m/s.
\begin{questions}
\item Calculer sa durée de propagation.
\item Exprimer cette durée en millisecondes.
\end{questions}
\end{exercice}\end{blocLivre}
\begin{blocLivre}{ex-14-4}\begin{exercice}[1]{Une distance sonore}
Un retard de 0,0500 s est mesuré pour un signal dans l’air à 340 m/s.
\begin{questions}
\item Calculer la distance parcourue.
\item La même durée donnerait-elle la même distance dans l’eau à 1500 m/s ?
\end{questions}
\end{exercice}\end{blocLivre}
\rubrique{S’entraîner}
\begin{blocLivre}{ex-14-5}\begin{exercice}[2]{Deux microphones}
Les microphones sont séparés de 1,02 m. Le retard lu vaut 3,00 ms.
\schema{retard-son}
\begin{questions}
\item Calculer la vitesse du son.
\item Pourquoi comparer le même détail des deux signaux ?
\end{questions}
\end{exercice}\end{blocLivre}
\begin{blocLivre}{ex-14-6}\begin{exercice}[2]{Un obstacle renvoie le son}
Un écho revient après 0,300 s ; $v=340$ m/s.
\schema{echo}
\begin{questions}
\item Déterminer la longueur du trajet total.
\item Déterminer la distance à l’obstacle.
\end{questions}
\end{exercice}\end{blocLivre}
\begin{blocLivre}{ex-14-7}\begin{exercice}[2]{Comparer des vitesses}
On donne 340 m/s dans l’air et 1500 m/s dans l’eau.
\begin{questions}
\item Calculer le rapport des vitesses.
\item Quel signal arrive le premier sur une même distance ?
\end{questions}
\end{exercice}\end{blocLivre}
\begin{blocLivre}{ex-14-8}\begin{exercice}[2]{Une matière qui reste}
Un son traverse un couloir de 10 m.
\begin{questions}
\item Les mêmes molécules d’air partent-elles de la source pour atteindre l’oreille ?
\item Que transmettent-elles de proche en proche ?
\end{questions}
\end{exercice}\end{blocLivre}
\rubrique{Raisonner et approfondir}
\begin{blocLivre}{ex-14-9}\begin{exercice}[3]{Position de la source}
Deux microphones sont distants de 1,0 m, mais la source n’est pas alignée avec eux.
\begin{questions}
\item Peut-on toujours prendre 1,0 m comme différence de trajets ?
\item Quelle disposition simple permet ce choix ?
\end{questions}
\end{exercice}\end{blocLivre}
\begin{blocLivre}{ex-14-10}\begin{exercice}[3]{Améliorer une résolution}
Le logiciel distingue des retards à 0,1 ms près. Une séparation trop petite donne un retard de 0,2 ms.
\begin{questions}
\item Pourquoi la mesure est-elle fragile ?
\item Que peut-on changer sans changer de logiciel ?
\end{questions}
\end{exercice}\end{blocLivre}
\begin{blocLivre}{ex-14-11}\begin{exercice}[3]{Un sonar}
Un sonar reçoit un retour 0,0800 s après émission. La vitesse du son dans l’eau vaut 1500 m/s.
\schema{echo}
\begin{questions}
\item Calculer la distance à l’objet.
\item Quelle erreur donnerait l’oubli de l’aller-retour ?
\end{questions}
\end{exercice}\end{blocLivre}
\begin{blocLivre}{ex-14-12}\begin{exercice}[3]{Un retard mal lu}
Un premier signal commence à 12,0 ms et le second à 15,0 ms. Les microphones sont séparés de 1,02 m.
\begin{questions}
\item Calculer le retard pertinent.
\item Calculer la vitesse et expliquer pourquoi 15,0 ms ne convient pas.
\end{questions}
\end{exercice}\end{blocLivre}

